\Titular*%
{historia}%
{A anoada orixe da teoría de nós}
{Víctor Díaz Díaz}
{Unha historia que mezcla o éter luminífero, unha caixa que fuma e a táboa periódica dos elementos.}

\begin{multicols}{2}

\subsection*{Introdución}

Dende atar os zapatos ata separar cables liados, os nós son un elemento moi común na vida cotiá de todo o mundo. Sen embargo, a pesar de ser tan comúns, o seu estudo rigoroso empregando as matemáticas tardou bastante tempo en poñerse en marcha, e as súas orixes están íntimamente relacionadas ca física: comezou co intento de dous científicos escoceses de aplicar estes obxectos á descrición dos elementos químicos.

Neste artigo imos repasar a curiosa orixe da teoría de nós moderna, que combina a física da luz e o éter, unha caixa chea de fume tóxico e a táboa periódica dos elementos.

\subsection*{O éter}

Para poñernos en contexto, temos que comezar por entender o que se coñecía da luz a finais do século XVII. Huygens postulou no seu tratado sobre a luz \cite{Huygens} que esta era unha onda lonxitudinal, que vibra na mesma dirección na que se propaga. Neste sentido, unha onda de luz comportaríase do mesmo xeito que unha onda mecánica nun fluido, similar ao son no aire. Esta perspectiva explicaba os fenómenos da refración e difración, pero non o de birefrixencia. Para intentar solucionar este problema, Newton propuxo que a luz estaba composta de corpúsculos \cite{Newton}. Esto lograba explicar a birrefrixencia, aínda que pola contra non conseguía recuperar a refración e a difración. Finalmente, a teoría ondulatoria impúxose: tratando a luz coma unha onda transversal, que vibra en direccións transversais á de propagación, podían explicarse os tres fenómenos de xeito satisfactorio.

Por outra banda, o consenso entre os físicos da época era que toda onda (xa sexa lonxitudinal ou transversal) necesita un medio físico para propagarse. Coñecíase que a luz era capaz de propagarse a unha velocidade altísima a través do aire, a auga e incluso o espazo, polo que todo o Universo debía estar cheo dun medio moi peculiar que servise de soporte ás ondas de luz. Este medio chamouse \textit{éter luminífero}, e sería esta substancia que vibraría para facer que a luz se propagase a través dela. Adicionalmente, este medio podería proporcionar un sistema de referencia privilexiado, solucionando outra cuestión importante na filosofía da física. 

Este é o contexto no que se desenvolveu o electromagnetismo clásico no século XIX, que orixinalmente non estaba conectado ca física da luz. Pola súa banda, Michael Faraday propuxo a existencia de \textit{liñas de forza} co obxectivo de explicar como poden as cargas eléctricas afectarse entre si. Faraday foi un científico eminentemente experimental, e non foi quen de desenvolver unha formulación matemática precisa das súas ideas. Este traballo foi feito posteriormente por James C. Maxwell.

Tratando de entender as liñas de forza, Maxwell notou que as ecuacións da electrostática no baleiro
\begin{equation}
	\nabla \cdot \textbf{E} = 0, ~~~~ \nabla \times \textbf{E} = 0,
\end{equation}
\noindent
teñen unha forma similar ás ecuacións do fluxo dun fluido incompresible sen viscosidade e vorticidade nula,
\begin{equation}
	\nabla \cdot \textbf{v} = 0, ~~~~ \nabla \times \textbf{v} = 0,
\end{equation}
\noindent
onde $\textbf{v}$ é o campo de velocidades. Esto deulle a entender que esas liñas de forza podían entenderse como pequenos tubos que transportan un fluido ideal. Este fluido sería o éter.

\subsection*{A caixa que fuma}

Agora imos relaxar a condición de vorticidade nula. Facendo esto, podemos introducir o concepto de \textit{liñas de vorticidade}. Unha liña de vorticidade dun campo $\textbf{v}$ ven dada por
\begin{equation}
	\omega = \nabla \times \textbf{v},
\end{equation}
\noindent
de xeito que é inmediato comprobar que a súa diverxencia é nula por definición,
\begin{equation}
	\nabla \cdot \omega = \nabla \cdot (\nabla \times \textbf{v}) = 0.
\end{equation}
\noindent
Esto implica que dita liña non pode ter fontes ou sumidoiros no interior do fluido, e únicamente pode ter os seus extremos na fronteira do fluido ou unidos entre eles, formando unha curva pechada. Hermann von Helmholtz demostrou a partir dos seus teoremas que unha liña de vorticidade nun fluido ideal é estable \cite{Helmholtz}. Esto implica que a liña pode deformarse, pero non pode romperse e disiparse no interior do fluido.  

Este resultado inspirou a Peter G. Tait para construir un mecanismo que producise tales liñas de vorticidade estables. O que deseñou foi unha caixa con dúas aberturas: nunha delas había un buraco circular e na outra unha membrana tersa. Esta caixa enchíase dalgunha substancia que producise fume e, ao golpear a membrana, este saía pola abertura circular formando aneis. Este truco xa era amplamente coñecido polos fumadores, pero Tait inventou a primeira caixa capaz de fumar.

\begin{figure}[H]
	\centering
	\includegraphics[width=\linewidth]{revistas/005/imaxes/caixa.png}
	\caption{A caixa construida por Tait. Extraído de \cite{Tait}.}
\end{figure}

A caixa producía aneis de fume que saían da súa boca e viaxaban polo aire, antes de desfacerse e disiparse por completo. Nin o fume do seu interior nin o aire no que viaxaban eran fluidos ideais (de feito, o fume era bastante tóxico), pero a construción servía para demostrar o efecto do que Helmholtz falaba. Tait mostroulle os seus resultados ao seu colega William Thompson (mellor coñecido como Lord Kelvin), quen quedou prendado do invento. Nos anos seguintes ambos traballarían en perfeccionar o aparello, e foi ao ver os aneis de fume viaxar polo aire do seu laboratorio cando Kelvin tivo unha idea revolucionaria: os distintos elementos químicos son liñas de vorticidade no éter \cite{Silver}.

\subsection*{A táboa periódica}

Uns anos antes, o científico inglés John Dalton introduciu o concepto do átomo na ciencia. O seu descubrimento xurdiu do estudo de compostos diferentes formados polos mesmos elementos, nos cales observou que as proporcións dos elementos eran sempre números enteiros pequenos \cite{Dalton}.

A conclusión de Dalton foi que a materia debía estar constituída por unidades indivisibles chamadas átomos. Sen embargo, a identidade destes átomos seguía sendo un misterio. O comportamento químico de cada elemento era moi diferente, o cal implicaba que cada tipo de átomo debía ter algunha propiedade intrínseca descoñecida que o diferenciase.

Foi precisamente esta clasificación a que Kelvin creu encontrar nos aneis de fume. A súa idea foi que un átomo é unha liña de vorticidade estabe no éter, correspondéndose cada elemento químico cun xeito diferente de formar un nó. Estes nós deben ter os extremos unidos e diferéncianse no sentido de que non poden ser deformados de xeito continuo (sen romper ou pegar ningún segmento) noutro nó. Dous nós que puidesen deformarse entre si, por moi complicado que fose o procedemento, serían equivalentes, dando lugar ao mesmo elemento. 

A táboa periódica dos elementos corresponderíase entón cunha táboa de nós non equivalentes, o que motivou a Tait a comezar a recoller nunha lista todos os nós diferentes. Este esforzo deu nacemento á teoría de nós, que se convertería nunha rama propia das matemáticas nos anos posteriores.

\begin{figure}[H]
	\centering
	\includegraphics[width=\linewidth]{revistas/005/imaxes/nos.png}
	\caption{Os sete nós máis sinxelos clasificados por Tait. Extraído de \cite{Silver}.}
\end{figure}

A idea tiña beleza e simplicidade, pero o seu fundamento científico estaba pouco definido. Sen embargo, eso non evitou que Kelvin se convencese de que o misterio dos átomos xa estaba resolto. E é que a pesar de ser un gran científico con ideas moi creativas, Lord Kelvin tiña o mal hábito de convencerse demasiado rápido das súas propias ideas. Xa en outras ocasións anteriores metera a pata coas súas conclusións, coma cando aseverou que a Terra non podía ter máis de 100 millóns de anos, indo en contra da evidencia xeolóxica e a recentemente publicada teoría da evolución por selección natural.

\subsection*{O fin do éter}

As malas noticias para Lord Kelvin chegarían en 1887, cando Albert A. Michaelson e Edward W. Morley intentaron detectar o éter co seu interferómetro, obtendo un resultado nulo \cite{MichaelsonMorley}. Aínda que este resultado non desbotaba por completo a existencia do éter, impoñía enormes restricións na súa posible natureza. 

O único xeito de reconciliar o resultado nulo ca existencia do éter era asumir que os obxectos se contraían na direción na que se movían a través do éter. Este fenómeno foi proposto independentemente por George F. FitzGerald \cite{FitzGerald} e Hendrik A. Lorentz \cite{Lorentz}, e posteriormente Henri Poincaré demostrou que era compatible co principio de relatividade \cite{Poincare} (i.e., as ecuacións describindo as leis físicas teñen a mesma forma en todos os sistema de referencia admisibles).

A revolución chegou en 1905 da man de Albert Einstein ca proposta dunha teoría derivada a partir de primeiros principios que ignoraba a existencia do éter, pero que conseguía explicar todos os fenómenos observables anteriormente atribuidos a el \cite{Einstein}. Esta nova teoría incluía o fenómeno de contracción de lonxitudes, e implicaba tanto a relatividade do espazo coma do tempo. 

\subsection*{Conclusións}

A historia da teoría de nós é curiosa e enrevesada, impulsada por Lord Kelvin e con orixes nunha intersección entre a física da luz e a química dos elementos. A proposta para intentar clasificar os elementos coma configuracións particulares do éter luminífero foi unha idea moi elegante, pero á que finalmente esa elegancia non a salvou de resultar sendo errónea. Lord Kelvin foi un dos máis grandes físicos da historia pero, coma xa se dixo, tiña a mala costume de quedar embelesado cas súas propias ideas demasiado pronto. Ademáis, daba un gran valor ao carácter estético das súas teorías. 

Esta práctica non é particularmente pouco común entre os físicos, especialmente en ámbitos teóricos. Dende que as matemáticas se empezaron a empregar para formular as leis físicas intentouse, ademáis, facelo de forma <<bela e elegante>>. Esto non é necesariamente contraproducente, e ten sido útil en algunhas ocasións. Pode axudar a descubrir conexións profundas e sutiles entre diferentes conceptos da física matemática que non son evidentes nun primeiro momento, pero debe aplicarse con coidado para non caer en sesgos subxectivos e convencerse de que unha teoría debe ser certa só pola súa beleza. Ao fin e ao cabo, a <<beleza>> nas matemáticas é un concepto fortemente subxectivo.

\printbibliography

\end{multicols}
